\section{Methodology}
\label{sec:method}

We study whether individual game experience transfers coordination skills to unrelated tasks with new partners.
Our experimental design pairs LLM agents from the same treatment condition and evaluates their ability to independently coordinate on tasks they have not seen during training.

\subsection{Experimental Conditions}
\label{sec:conditions}

We compare five conditions, summarized in \tabref{tab:conditions}.

\begin{table}[t]
\centering
\caption{Experimental conditions. Each condition provides a different ``game experience'' via the agent's system prompt. The AI-Designed game (\syncminds) combines all four coordination skills.}
\label{tab:conditions}
\small
\begin{tabular}{@{}llll@{}}
\toprule
\textbf{Condition} & \textbf{Game} & \textbf{Skills Trained} & \textbf{Purpose} \\
\midrule
\nogame & None & None & Baseline \\
\trivia & \brainquest & None (irrelevant) & Active control \\
\focalpt & \mindmatch & Focal point identification & Single-skill \\
\convention & \signalcraft & Convention formation & Single-skill \\
\aidesigned & \syncminds & All four skills & Multi-skill \\
\bottomrule
\end{tabular}
\end{table}

\para{Baseline (\nogame).}
Agents receive no game context, establishing raw LLM coordination ability.

\para{Active control (\trivia).}
Agents are told they have played \brainquest, a trivia game with no coordination component.
This controls for the effect of having \emph{any} game experience in the system prompt.

\para{Focal point game (\focalpt).}
Agents are given experience with \mindmatch, a game that trains identification of culturally salient choices---the ``obvious'' answer that a stranger would also pick.

\para{Convention game (\convention).}
Agents are given experience with \signalcraft, a game that trains formation of arbitrary shared conventions through repeated interaction.

\para{AI-designed multi-skill game (\aidesigned).}
Agents are given experience with \syncminds, a game designed by \gptfour to maximize coordination transfer.
\syncminds combines four coordination skills: focal point identification, perspective-taking (\tom), role differentiation, and convention formation.
The key design insight is that effective coordination requires both \emph{convergence} (matching your partner) and \emph{differentiation} (complementing your partner), and the game explicitly trains both.

\subsection{Evaluation Tasks}
\label{sec:tasks}

We evaluate coordination transfer on 31 tasks across three domains.

\para{Focal point tasks (22 tasks).}
Two agents independently choose a response to an open-ended prompt (\eg ``pick a number'' or ``pick a city'').
Coordination succeeds when both agents choose the same answer.
We include 10 \emph{standard} tasks with strong cultural focal points (\eg pick-a-color, pick-a-day) and 12 \emph{hard} tasks with larger option spaces or weaker salience (\eg pick-a-number from 1--1000, pick-an-emoji, pick-a-movie-genre).

\para{Role assignment tasks (9 tasks).}
Two or three agents must independently choose \emph{different} roles from a set of options (\eg cook vs.\ server, or three-way job assignment).
Coordination succeeds when agents adopt complementary, non-overlapping roles.
We include 5 standard binary-choice tasks and 4 hard tasks with 3-way assignment, resource allocation, and timing coordination.

\para{Convention formation (1 task).}
A Lewis signaling game~\citep{lipowska2018emergence}: two agents must develop a shared mapping between 3 objects and 3 signals over 10 rounds, receiving only success/failure feedback.
We measure convergence rate, average rounds to convergence, and final accuracy.

\subsection{Agent Setup}
\label{sec:agents}

We use \gptfour as the base model for all agents, following \citet{riedl2026emergent}'s validation of LLM agents as coordination proxies.

\para{Simulation of game experience.}
We encode ``having played a game'' by including a description of the game experience in each agent's system prompt.
This is analogous to how game experience would be encoded in a human player's memory and decision-making heuristics---the agent ``remembers'' the strategies and skills learned during gameplay.

\para{Independence.}
Within each pair, agents receive independent random seeds and have no access to each other's responses.
All coordination must be tacit: agents cannot communicate during the evaluation tasks.

\para{Hyperparameters.}
We set temperature $= 0.9$ (high enough for response diversity while maintaining coherent responses), max tokens $= 50$, and random seed $= 42$.
Each condition--task combination is evaluated with $N{=}50$ independent agent pairs, yielding approximately 12,000 total API calls across all conditions and tasks.

\para{Response normalization.}
Agent responses are normalized for comparison: case folding, whitespace trimming, and common synonym resolution (\eg ``New York'' $\to$ ``new york'', ``NYC'' $\to$ ``new york'').

\subsection{Evaluation Metrics}
\label{sec:metrics}

\para{Match rate.}
For focal point tasks, the fraction of pairs that choose the same (normalized) response.

\para{Complementarity score.}
For role assignment tasks, the fraction of pairs that adopt non-overlapping roles.

\para{Convention speed.}
For the signaling game, the average number of rounds to reach consistent agreement, along with convergence rate and final accuracy.

\para{Coordination improvement ratio.}
The primary metric: $\text{ratio} = \text{treatment match rate} / \text{baseline match rate}$.
A ratio of $2.0\times$ means coordination doubled relative to the untrained baseline.
